\begin{lemma}
Let \(I\subseteq[n]\) be a fixed set with
\[
|I|=k=\noderef{TwoBiteNaturalIndependenceScale}(1+\varepsilon,n).
\]
Assume \(\varepsilon_1\ge0\),
\(\noderef{ParameterHierarchyEventualComparisons}
(\varepsilon,\varepsilon_1,\varepsilon_2,n_0)\), \(n>n_0\), and the rounded
sample-parameter identities
\[
m=\noderef{TwoBiteNaturalReducedVertexCount}(n),\qquad
p=\noderef{TwoBiteEdgeProbability}\!\left(\frac12,n\right).
\]
If a sample \(\omega\) satisfies
\[
\noderef{TwoBiteRegularityEvent}(k,\varepsilon,\varepsilon_1,\varepsilon_2,\omega)
\quad\text{and}\quad
\noderef{TwoBiteProjectionSizeEvent}(I,k,\ell_R,\ell_B,\omega),
\]
then the staged-revelation open-pair set satisfies
\[
\noderef{ProjectionOpenPairFunction}(\ell_R,\ell_B,k,p,n)
  -10\varepsilon_1 k^2
\le
\left|\noderef{TwoBiteStagedOpenPairs}(\omega,\varepsilon,I)\right|.
\]
\end{lemma}
\begin{proof}
Keep the regularity hypothesis and the projection-size hypothesis separate.
First unfold the projection-size event
\(\noderef{TwoBiteProjectionSizeEvent}(I,k,\ell_R,\ell_B,\omega)\).  By its
definition this is a threefold conjunction, and its three projections are the
natural-number equalities
\[
h_I:\ |I|=k,\qquad
h_R:\ |I.\operatorname{image}\pi_R^\omega|=\ell_R,\qquad
h_B:\ |I.\operatorname{image}\pi_B^\omega|=\ell_B, \tag{1}
\]
where \(\pi_R^\omega=\noderef{TwoBiteRedProjection}(\omega)\) and
\(\pi_B^\omega=\noderef{TwoBiteBlueProjection}(\omega)\).  The equality
\(h_I\) is the same scale equality as the separately stated active hypothesis
\(|I|=k\); either copy may be used for rewriting, and below we use the one
coming directly from the projection-size event.  Casting (1) to \(\mathbb R\)
gives
\[
h_I^{\mathbb R}:\ (|I|:\mathbb R)=(k:\mathbb R),\quad
h_R^{\mathbb R}:\ (|I.\operatorname{image}\pi_R^\omega|:\mathbb R)
   =(\ell_R:\mathbb R),\quad
h_B^{\mathbb R}:\ (|I.\operatorname{image}\pi_B^\omega|:\mathbb R)
   =(\ell_B:\mathbb R). \tag{2}
\]

The child \(\noderef{RevealingLeavesOpenPairs}\) does not take the external
scale \(k\) as an input; it uses the scale \(|I|\).  Its independence-scale
hypothesis is therefore
\[
|I|=\noderef{TwoBiteNaturalIndependenceScale}(1+\varepsilon,n).
\tag{3}
\]
This follows by transitivity from \(h_I:|I|=k\) and the active scale
hypothesis
\[
k=\noderef{TwoBiteNaturalIndependenceScale}(1+\varepsilon,n).
\tag{4}
\]
No regularity information is used in this step.

Now unfold the regularity event
\(\noderef{TwoBiteRegularityEvent}(k,\varepsilon,\varepsilon_1,\varepsilon_2,
\omega)\).  Its first, second, and third conjuncts, in the order of the Lean
definition, are exactly
\[
\noderef{FiberAndDegreeEvent}(\omega,\varepsilon_2), \tag{5a}
\]
\[
\noderef{TwoBiteMediumClosedPairsEvent}(k,\varepsilon,\varepsilon_1,\omega),
\tag{5b}
\]
\[
\noderef{TwoBiteSmallClosedPairsEvent}(k,\varepsilon,\varepsilon_1,\omega).
\tag{5c}
\]
The first conjunct (5a) is already the fiber-degree hypothesis required by
\(\noderef{RevealingLeavesOpenPairs}\).  The two closed-pair conjuncts must be
transported from the displayed scale \(k\) to the child theorem's scale
\(|I|\).  Since \(h_I:|I|=k\), equality elimination for the family
\[
s\longmapsto
\noderef{TwoBiteMediumClosedPairsEvent}(s,\varepsilon,\varepsilon_1,\omega)
\]
rewrites the natural-number index \(k\) in (5b) to \(|I|\).  Equivalently,
substitute \(k=|I|\) in the defining predicate
\[
\forall J,\ |J|=k\to
\text{the closed-pair bound at scale }k,
\]
including both the premise \(|J|=k\) and the real scale appearing in the
closed-pair bound.  This gives
\[
\noderef{TwoBiteMediumClosedPairsEvent}(|I|,\varepsilon,\varepsilon_1,\omega).
\tag{6b}
\]
The same equality transport for the analogous family
\[
s\longmapsto
\noderef{TwoBiteSmallClosedPairsEvent}(s,\varepsilon,\varepsilon_1,\omega)
\]
applied to (5c) gives
\[
\noderef{TwoBiteSmallClosedPairsEvent}(|I|,\varepsilon,\varepsilon_1,\omega).
\tag{6c}
\]
These transports change only the explicit index of the event; the sample,
the parameters \(\varepsilon,\varepsilon_1\), and the closed-pair event
families themselves are unchanged.

All remaining hypotheses of \(\noderef{RevealingLeavesOpenPairs}\) are passed
verbatim from the active assumptions:
\[
0\le\varepsilon_1,\qquad
\noderef{ParameterHierarchyEventualComparisons}
(\varepsilon,\varepsilon_1,\varepsilon_2,n_0),\qquad n_0<n,
\tag{7}
\]
and the rounded parameter identities
\[
m=\noderef{TwoBiteNaturalReducedVertexCount}(n),\qquad
p=\noderef{TwoBiteEdgeProbability}\!\left(\frac12,n\right).
\tag{8}
\]
Applying \(\noderef{RevealingLeavesOpenPairs}\) to the same
\(\omega,I,\varepsilon,\varepsilon_1,\varepsilon_2,n_0,m,p\) now matches its
hypotheses term by term: (7) supplies \(0\le\varepsilon_1\), the hierarchy
package, and \(n_0<n\); (8) supplies the rounded \(m\) and \(p\) identities;
(3) supplies the child scale equality \( |I|=\noderef{TwoBiteNaturalIndependenceScale}(1+\varepsilon,n)\);
(5a) supplies the fiber-degree event; and (6b),(6c) supply the medium and
small closed-pair events at the child index \(k:=|I|\).  The child conclusion
therefore yields
\[
\begin{aligned}
&\noderef{ProjectionOpenPairFunction}
 \left((|I.\operatorname{image}\pi_R^\omega|:\mathbb R),
       (|I.\operatorname{image}\pi_B^\omega|:\mathbb R),
       (|I|:\mathbb R),p,(n:\mathbb R)\right)
 -10\varepsilon_1 (|I|:\mathbb R)^2  \\
&\hspace{5em}\le
\left(\left|
\noderef{TwoBiteStagedOpenPairs}(\omega,\varepsilon,I)
\right|:\mathbb R\right).
\end{aligned}
\tag{9}
\]

It remains only to rewrite the left side of (9) from the child theorem's
internal projection sizes to the active statement's named integers
\(\ell_R,\ell_B,k\).  The function
\(\noderef{ProjectionOpenPairFunction}\) is an ordinary five-argument real
function, so congruence allows substitutions in any of its first three
arguments while leaving the final two arguments \(p\) and \((n:\mathbb R)\)
fixed.  Substitute \(h_R^{\mathbb R}\) into the first argument, substitute
\(h_B^{\mathbb R}\) into the second argument, and substitute
\(h_I^{\mathbb R}\) into the third argument.  The same equality
\(h_I^{\mathbb R}\) also rewrites the penalty term:
\[
(|I|:\mathbb R)^2=(k:\mathbb R)^2,
\quad\text{hence}\quad
10\varepsilon_1(|I|:\mathbb R)^2
  =10\varepsilon_1(k:\mathbb R)^2.
\tag{10}
\]
These substitutions are made inside the whole left-hand side
\[
x\longmapsto
\noderef{ProjectionOpenPairFunction}(\cdots)-10\varepsilon_1(\cdots)^2,
\]
so the subtraction is rewritten by the same equality congruence.  The
arguments \(p\) and \((n:\mathbb R)\), and the right-hand side
\(\left(|\noderef{TwoBiteStagedOpenPairs}(\omega,\varepsilon,I)|:\mathbb R\right)\),
are not changed by these substitutions.  After the three argument rewrites and
the penalty rewrite, (9) is exactly
\[
\noderef{ProjectionOpenPairFunction}
  ((\ell_R:\mathbb R),(\ell_B:\mathbb R),(k:\mathbb R),p,(n:\mathbb R))
  -10\varepsilon_1 (k:\mathbb R)^2
\le
\left(\left|
\noderef{TwoBiteStagedOpenPairs}(\omega,\varepsilon,I)
\right|:\mathbb R\right),
\]
which is the asserted open-pair lower bound.
\end{proof}
